\chapter{Serii oarecare. Aproximații}


\section{Convergență absolută și semiconvergență}

Dacă lucrăm cu serii care pot avea și termeni negativi, studiul convergenței
trebuie făcut mai atent. Astfel, avem nevoie de următoarele:

\begin{definition}\label{def:serie-alt} \index{serii!alternate}
  O serie $ \sum_n x_n $ se numește \emph{alternată} dacă produsul $ x_n \cdot x_{n+1} < 0 $,
  pentru orice indice $ n \in \NN $.
\end{definition}

Pentru asemenea serii, avem la dispoziție un singur criteriu, anume:

\begin{theorem}[Criteriul lui Leibniz]\label{thm:leibniz} \index{serii!criteriu!Leibniz}
  Fie $ \sum_n (-1)^n x_n $ o serie alternată.

  Dacă șirul $ (x_n) $ este descrescător și converge către zero, atunci seria
  este convergentă.
\end{theorem}

Pentru serii alternate, avem următoarele noțiuni suplimentare:

\begin{definition}\label{def:ac-sc}\index{serii!absolut convergente} \index{serii!semiconvergente}
  O serie $ \sum_n x_n $ se numește \emph{absolut convergentă} dacă seria modulelor
  $ \sum_n |x_n| $ este convergentă și \emph{semiconvergentă} dacă are seria
  modulelor divergentă.
\end{definition}

Se poate constata imediat că \emph{orice serie absolut convergentă este convergentă},
deoarece pentru orice număr natural $ x $, avem $ x \leq |x| $. În plus, dacă
seria inițială este divergentă, atunci și seria modulelor va fi divergentă, din
același motiv.

Pentru studiul convergenței seriilor generale, adică avînd și termeni pozitivi,
și negativi, avem un criteriu important.

\begin{theorem}[Criteriul Abel-Dirichlet]\label{thm:crit-abel} \index{serii!criteriu!Abel-Dirichlet}
  Presupunem că seria $ \sum_n x_n $ se poate scrie sub forma $ \sum_n \alpha_n y_n $, cu
  $ (\alpha_n) $ un șir monoton și mărginit. Dacă seria $ \sum_n y_n $ este convergentă,
  atunci și seria inițială este convergentă.

  \textbf{Alternativ}, uneori criteriul este formulat astfel:
  Dacă $ (\alpha_n) $ este un șir monoton care tinde către zero, iar
  șirul cu termenul general $ Y_n = y_1 + \dots + y_n $ este mărginit, atunci seria
  inițială este convergentă.
\end{theorem}

De exemplu, să studiem seria $ \ds\sum_n \dfrac{(-1)^n}{n} $. Fiind o serie
alternată, constatăm:
\begin{itemize}
\item Seria modulelor este $ \ds\sum_n \dfrac{(-1)^n}{n} $, care este seria
  armonică, deci divergentă;
\item Pentru seria inițială, putem aplica criteriul lui Leibniz. Șirul $ x_n = \dfrac{1}{n} $
  este evident descrescător către zero, deci seria este convergentă.
\end{itemize}

Concluzia este că seria este \emph{semiconvergentă}.

%%%%%%%%%%%%%%%%%%%%%%%%%%%%%%%%%%%%%%%%%%%%%%%%%%%%%%%%%%%%%%%%%%%%%%

\section{Aproximarea sumelor seriilor convergente}

Putem calcula numeric sumele unor serii convergente, cu aproximații oricît de bune.
Următoarele două rezultate sînt fundamentale:

\begin{theorem}[Aproximarea sumelor seriilor cu termeni pozitivi]\label{thm:aprox-poz}
  Fie $ x_n \geq 0 $ și $ k \geq 0 $, astfel încît să avem:
  \[
    \dfrac{x_{n+1}}{x_n} < k < 1, \quad \forall n \in \NN.
  \]

  Dacă $ S $ este suma seriei convergente $ \sum_n x_n $, iar $ s_n $ este suma
  primilor $ n $ termeni, atunci avem aproximația:
  \[
    |S - s_n| < \dfrac{k}{1 - k} x_n.
  \]
\end{theorem}

\begin{theorem}[Aproximarea sumelor seriilor alternate]\label{thm:aprox-alt}
  Fie $ \sum_n (-1)^n x_n $ o serie alternată, convergentă și fie $ S $ suma sa.

  Dacă $ S_n $ este suma primilor $ n $ termeni, atunci avem:
  \[
    | S - s_n | \leq x_{n + 1}.
  \]
  Cu alte cuvinte, eroarea aproximației are ordinul de mărime al primului
  termen neglijat.
\end{theorem}

%%%%%%%%%%%%%%%%%%%%%%%%%%%%%%%%%%%%%%%%%%%%%%%%%%%%%%%%%%%%%%%%%%%%%%

\section{Exerciții}

1. Studiați natura (AC/SC/D) următoarelor serii, definite de șirul $ x_n $:
\begin{enumerate}[(a)]
\item $ x_n = (-1)^n $;
\item $ x_n = \dfrac{(-1)^n}{n} $;
\item $ x_n = \dfrac{1}{n + i} $;
\item $ x_n = \dfrac{1}{(n + i)\sqrt{n}} $;
\item $ x_n = (-1)^n \dfrac{\log_a n}{n}, a > 1 $;
\item $ x_n = (-1)^{n+1} \dfrac{2n + 1}{3^n} $;
\item $ x_n = (-1)^{n+1} \dfrac{1}{\sqrt{n(n+1)}} $.
\end{enumerate}

\vspace{1cm}

2. Să se aproximeze cu o eroare mai mică decît $ \varepsilon $ sumele seriilor
definite de șirul $ x_n $:
\begin{enumerate}[(a)]
\item $ x_n = \dfrac{(-1)^n}{n!}, \varepsilon = 10^{-3} $ ($ n = 7 $);
\item $ x_n = \dfrac{(-1)^n}{n^3 \sqrt{n}}, \varepsilon  = 10^{-2} $ ($ n = 4 $);
\item $ x_n = \dfrac{1}{n! 2^n}, \varepsilon = 10^{-4} $;
\item $ x_n = \dfrac{(-1)^n}{n^3 \sqrt{n}}, \varepsilon = 10^{-2} $.
\end{enumerate}

\vspace{1cm}

3. Să se demonstreze convergența și să se determine sumele seriilor cu
termenul general dat de:
\begin{enumerate}[(a)]
\item $ x_n = \dfrac{1}{(\alpha + n)(\alpha + n + 1)} $;
\item $ x_n = \dfrac{1}{16n^2 - 8n - 3} $;
\item $ x_n = \ln \dfrac{n+1}{n} $.
\end{enumerate}


%%% Local Variables:
%%% mode: latex
%%% TeX-master: "../bcd-18-19"
%%% End:
